\setcounter{framenumber}{0}
%22

\begin{frame}
  \titlepage
\end{frame}

\begin{frame}{Learning Goals}
By the end of this lecture, students should be able to:
\begin{itemize}
    \item understand the Law of Large Numbers;
    \item distinguish convergence in probability from exact equality;
    \item apply the Weak Law of Large Numbers;
    \item understand Monte Carlo integration;
    \item define the empirical CDF;
    \item understand the Central Limit Theorem;
    \item standardize sums and averages;
    \item apply Normal approximations to Binomial and Poisson distributions.
\end{itemize}
\end{frame}

\lecturepart{Why Limit Theorems Matter}

\begin{frame}{Order Emerging from Randomness}
Individual random outcomes are unpredictable.

\vspace{0.25cm}

But large collections of random variables often display remarkable regularity.

\vspace{0.25cm}

Examples:
\begin{itemize}
    \item sample averages stabilize;
    \item proportions settle down;
    \item histograms become bell-shaped;
    \item randomness averages out.
\end{itemize}

\vspace{0.25cm}

\begin{keybox}{Main message}
Large-scale randomness becomes highly structured.
\end{keybox}
\end{frame}

\begin{frame}{Two Fundamental Theorems}
The two most important limit theorems are:

\vspace{0.25cm}

\begin{enumerate}
    \item Law of Large Numbers (LLN)

    \[
    \text{averages converge to expectations;}
    \]

    \item Central Limit Theorem (CLT)

    \[
    \text{standardized sums become approximately Normal.}
    \]
\end{enumerate}

\vspace{0.25cm}

Together, these form the foundation of modern statistics.
\end{frame}

\lecturepart{Law of Large Numbers (LLN)}


\begin{frame}{Weak Law of Large Numbers}

\begin{keybox}{Theorem (WLLN)}
Let $X_1,X_2,\dots$ be i.i.d.\ random variables with finite mean $\mu$ (and finite variance $\sigma^2)$


Then
\[
\forall  \ \varepsilon>0, \quad  \lim_{n\to\infty}
\Prob\!\left(
|\bar X_n-\mu|>\varepsilon
\right)
=0, \quad 
\text{or equivalently,} 
\quad   \lim_{n\to\infty}
\Prob\!\left(
|\bar X_n-\mu|\le \varepsilon
\right)
=1.
\]
 
\end{keybox}



\begin{keybox}{Interpretation}
As the sample size grows, the sample mean $\bar X_n$
becomes arbitrarily close to the population mean $\mu$
with probability approaching \(1\).
\end{keybox}

\textbf{Notation (Convergence in Probability):} This phenomena is also known as \textit{convergence in probability}, and is denoted by
\begin{equation*}
    \bar X_n \xrightarrow[]{P}\mu.
\end{equation*}

\end{frame}






\begin{frame}{Proof of WLLN Using Chebyshev's Inequality}
Recall:
\[
\mathbb{E}(\bar X_n)=\mu,
\]
and
\[
\operatorname{Var}(\bar X_n)
=
\frac{\sigma^2}{n}.
\]

Applying Chebyshev:
\[
\Prob(|\bar X_n-\mu|>\varepsilon)
\le
\frac{
\operatorname{Var}(\bar X_n)
}{
\varepsilon^2
}.
\]

Thus:
\[
\Prob(|\bar X_n-\mu|>\varepsilon)
\le
\frac{\sigma^2}{n\varepsilon^2}.
\]

As
\[
n\to\infty,
\]
the right-hand side goes to
\[
0.
\]
\end{frame}
\begin{frame}{Coin Toss Example}

Suppose:
\[
X_i=
\begin{cases}
1,&\text{Heads},\\
0,&\text{Tails}.
\end{cases}
\]

Then $\mathbb{E}(X_i)=p.$ The sample mean is:
\[
\bar X_n
=
\frac{1}{n}\sum_{i=1}^n X_i
=
\text{proportion of Heads}.
\]
The LLN says $X_n
\overset{P}{\longrightarrow}
p.$
Thus, observed proportions stabilize near the true probability \(p\).

\end{frame}

\begin{frame}{LLN and the Relative-Frequency Interpretation of Probability}

\begin{keybox}{}
The Law of Large Numbers gives a mathematical justification for the
\textbf{long-run relative-frequency interpretation} of probability:
if an experiment is repeated many times, the proportion of trials on which
an event occurs approaches its probability.
\end{keybox}

Let \(A\subseteq \Sample\) be an event, and suppose the experiment is repeated
independently under the same conditions. Define
\[
X_i=
\begin{cases}
1, & \text{if } A \text{ occurs on trial } i,\\
0, & \text{otherwise}.
\end{cases}
\]

Then
\[
X_i\sim \operatorname{Bernoulli}(\Prob(A)),
\qquad
\E(X_i)=\Prob(A).
\]

Moreover,
\[
\bar X_n
=
\frac{1}{n}\sum_{i=1}^n X_i
=
\frac{\#\{\text{times }A\text{ occurs in the first }n\text{ trials}\}}{n},
\]
so \(\bar X_n\) is exactly the \textbf{relative frequency} of \(A\). By the Law of Large Numbers, $\bar X_n
\overset{P}{\longrightarrow}
\Prob(A).$ Thus,
\[
\text{relative frequency of }A
\quad
\xrightarrow[\text{number of trials}\to\infty]{P}
\quad
\text{theoretical probability }\Prob(A).
\]

\end{frame}



\begin{frame}{Strong Law of Large Numbers}
\begin{keybox}{Strong Law}
Under the same assumptions of WLLN, we have

\begin{equation*}
    \Prob({\lim_{n\to \infty} \bar X_n = \mu}) =1.
\end{equation*}
Equivalently, 
\begin{equation*}
    \Prob(\{ s\in \Sample : \lim_{n\to \infty}X(s)=\mu \})=1
\end{equation*}

\end{keybox}

\textbf{Notation (Almost-Sure Convergence)}: This phenomenon is also known as \textit{almost-sure convergence} or \textit{almost-everywhere convergence}, and is denoted by
\begin{equation*}
    \bar X_n \xrightarrow[]{a.s.}\mu, \quad or \quad \bar X_n \xrightarrow[]{a.e.}\mu
\end{equation*}

\textbf{Remark}: Almost-sure convergence implies (i.e., is stronger than) convergence in probability. Thus, SLLN is stronger than WLLN. The proof requires graduate-level background of Probability (Borel-Canteli Lemma).


\end{frame}

\lecturepart{LLN Application: Monte Carlo Integration}
\begin{frame}{Monte Carlo Integration}

Suppose we want to compute
\[
\int_0^1 g(x)\,dx.
\]

Let
$U_1,\dots,U_n
\overset{\text{i.i.d.}}{\sim}
\operatorname{Unif}(0,1)$.

Then
\[
\mathbb{E}(g(U_i))
=
\int_0^1 g(x)\,dx.
\]

Hence,
\[
\mathbb{E}[\frac1n
\sum_{j=1}^n g(U_j)]= \int_0^1 g(x)\,dx.
\]
By the WLLN,
\[
\boxed{
\frac1n
\sum_{j=1}^n g(U_j)
\xrightarrow[]{P}
\int_0^1 g(x)\,dx, \quad \text{and thus},\quad   \frac1n
\sum_{j=1}^n g(U_j)
\approx
\int_0^1 g(x)\,dx
\quad
\text{for large } n.
}
\]

This is called Monte Carlo estimation of integrals. Draw $n$ random points on the interval, evaluate the function $g$ on them, and then average. The larger the $n$ the closer the average to the true integral.

\end{frame}



\lecturepart{LLN Application: Empirical Distribution Function}

\begin{frame}{Empirical CDF}
Given observations:
\[
X_1,\dots,X_n \sim F
\]
define the empirical CDF:
\[
{
\hat F_n(x)
=
\frac1n
\sum_{j=1}^n
I(X_j\le x).
}
\]

Interpretation:
\[
\hat F_n(x)
=
\text{fraction of observations $\leq$ }x.
\]


For fixed $x,$ the indicators $I(X_j\le x)$ are i.i.d. Bernoulli random variables with mean $F(x)=\Prob(X\le x).$ Thus, by the LLN:
\[
{
\hat F_n(x)\to F(x).
}
\]

The empirical CDF approximates the true CDF.
\end{frame}




\lecturepart{Central Limit Theorem}

\begin{frame}{Central Limit Theorem}

\begin{keybox}{Theorem (CLT)}
Let $X_1,X_2,\dots$ be i.i.d.\ random variables with
$\mathbb{E}(X_i)=\mu$ and $\Var(X_i)=\sigma^2<\infty$.

Then for every $z\in \mathbb{R}$,
\[
\forall z\in \mathbb{R},\ \lim_{n\to\infty}
\Prob\!\left(
\frac{
(\bar X_n-\mu)
}{
\sigma /\sqrt n
}
\le z
\right)
=
\Phi(z),
\]
\end{keybox}


\begin{keybox}{Interpretation}
We previously learned that $\mathbb{E}[\bar X_n] = \mu$ and $\operatorname{Var}[\bar X_n]=\frac{\sigma^2}{n}$, which means that the fraction $\frac{\bar X_n-\mu}{\sigma /\sqrt{n}}$ is the standardized version of $\bar X_n$. CLT says that for large \(n\), the sample mean
$\bar X_n$
is approximately normal: $\bar X_n
\approx
N\!\left(
\mu,
\frac{\sigma^2}{n}
\right).$
This is especially important because we do not assume any distribution for $X_1,\dots X_n$ yet the sample mean follows a normal distribution for large $n.$ This offers that Normality emerges universally from averaging.
\end{keybox}

\textbf{Notation (Convergence in Distribution):} This phenomenon is also known as \textit{convergence in distribution}, and is denoted by $\frac{
 \bar X_n-\mu
}{
\sigma/\sqrt n
}
\xrightarrow[]{d}
N(0,1).$
\end{frame}



\begin{frame}{CLT Application: Normal's Approximation to Binomial}
We have previously seen that a binomial r.v. can be written as a sum of $n$ i.i.d. Bernouli r.v.s. Let $X\sim \operatorname{Bin}(n,p).
$, and $X
=
X_1+\cdots+X_n,$ where $X_1,\dots, X_n \sim Ber(p)$.


Since:
\[
\mu=\mathbb{E}[X_i]= p,
\qquad
\sigma^2=\operatorname{Var}[X_i]= p(1-p),
\]
the CLT gives:
\[
\boxed{
\frac{\bar X_n-\mathbb{E}[\bar X_n]}{\sqrt{Var}[\bar X_n]}=\frac{X/n-p}{\sqrt{p(1-p)/n}}
=
\frac{
X-np
}{
\sqrt{np(1-p)}
}
\approx
N(0,1).
}
\]
\end{frame}

\begin{frame}{Continuity Correction}
A subtle point here is that Binomial distributions are discrete, while Normal distributions are continuous.

\vspace{0.25cm}

A better approximation uses a continuity correction:
\[
\Prob(X\le k)
\approx
\Prob\left(
Y\le k+\frac12
\right),
\]
where
\[
Y\sim N(np,np(1-p)).
\]

\vspace{0.25cm}

\begin{keybox}{Idea}
The correction accounts for the width of discrete bins.
\end{keybox}
\end{frame}

\begin{frame}{Poisson to Normal}
Suppose:
\[
X\sim \operatorname{Pois}(\lambda).
\]

For large
\[
\lambda,
\]
\[
{
\frac{
X-\lambda
}{
\sqrt{\lambda}
}
\approx
N(0,1).
}
\]

\vspace{0.25cm}

Equivalently:
\[
X
\approx
N(\lambda,\lambda).
\]
\end{frame}

\begin{frame}{Why the CLT Is Remarkable}
The original distribution can be:
\begin{itemize}
    \item discrete;
    \item continuous;
    \item skewed;
    \item highly non-Normal.
\end{itemize}

\vspace{0.25cm}

Yet averages still become approximately Normal.

\vspace{0.25cm}

\begin{keybox}{Universality}
The CLT explains why Normal distributions appear everywhere.
\end{keybox}
\end{frame}


\begin{frame}{Summary}
\begin{itemize}
    \item The LLN says averages stabilize.

    \item Monte Carlo integration is justified by the LLN.

    \item The empirical CDF approximates the true CDF.

    \item The CLT says standardized sums/averages become approximately Normal.

    \item Standard errors decrease like
    \[
    \frac1{\sqrt n}.
    \]

    \item Normal approximations apply to Binomial and Poisson distributions.
\end{itemize}

\vspace{0.25cm}

\begin{keybox}{Next lecture}
Chi-Square and Student $t$ distributions.
\end{keybox}
\end{frame}